\ifdefined\gkver\else\def\gkver{student}\fi
\documentclass[\gkver]{gaokaozhenti}
\begin{document}

\chapter{2020年1月浙江}

\begin{enumerate}
\item
%% number: 1
%% paperName: 2020年1月浙江选考第 1 题 3 分
%% typeId: 1
%% type: 单选题
%% chapter: 运动和力
%% point: 力学单位制
%% method: 无
%% score: 3
%% degree: 850
%% duplicateId: 0
%% body:
以下物理量为矢量，且单位是国际单位制基本单位的是\xzanswer{B}
\fourchoices[answer=B]
{电流、A}
{位移、m}
{功、J}
{磁感应强度、T}
%% answer: B
%% memo:
\memoanswer{
电流是标量，其单位为国际单位制中的基本单位 A，故 A 错误；位移是矢量，其单位为国际单位制中的基本单位 m，故 B 正确；功是标量，其单位不是国际单位制中的基本单位，故 C 错误；磁感应强度是矢量，其单位不是国际单位制中的基本单位，故 D 错误。
}

\item
%% number: 2
%% paperName: 2020年1月浙江选考第 2 题 3 分
%% typeId: 1
%% type: 单选题
%% chapter: 运动和力
%% point: 牛顿第三定律
%% method: 无
%% score: 3
%% degree: 850
%% duplicateId: 0
%% body:
如图所示，一对父子掰腕子，父亲让儿子获胜。若父亲对儿子的力记为 $F_{1}$，儿子对父亲的力记为 $F_{2}$，则\xzanswer{B}
\onepicture[w=5cm]{figs/02.png}
\fourchoices[answer=B]
{$F_{1} > F_{2}$}
{$F_{1}$ 和 $F_{2}$ 大小相等}
{$F_{1}$ 先于 $F_{2}$ 产生}
{$F_{1}$ 后于 $F_{2}$ 产生}
%% answer: B
%% memo:
\memoanswer{
由牛顿第三定律知，$F_{1}$、$F_{2}$ 大小相等，且同时产生，故 B 正确，A、C、D 错误。
}

\item
%% number: 3
%% paperName: 2020年1月浙江选考第 3 题 3 分
%% typeId: 1
%% type: 单选题
%% chapter: 直线运动
%% point: 参考系
%% method: 无
%% score: 3
%% degree: 850
%% duplicateId: 0
%% body:
如图所示，新中国成立 70 周年阅兵仪式上，国产武装直升机排列并保持“70”字样编队从天安门上空整齐飞过。甲、乙分别是编队中的两架直升机，则\xzanswer{D}
\onepicture[w=7cm]{figs/03.jpg}
\fourchoices[answer=D]
{以甲为参考系，乙是运动的}
{以乙为参考系，甲是运动的}
{以甲为参考系，坐在观众席上的观众都是静止的}
{以乙为参考系，“70”字样编队中所有直升机都是静止的}
%% answer: D
%% memo:
\memoanswer{
直升机编队保持“70”字样飞行，则以甲为参考系，甲、乙间没有位置的改变，故乙是相对静止的；同理以乙为参考系，甲是静止的，故 A、B 错误。以甲为参考系，坐在观众席上的观众都是运动的，故 C 错误。以乙为参考系，“70”字样编队中所有直升机都是静止的，故 D 正确。
}

\item
%% number: 4
%% paperName: 2020年1月浙江选考第 4 题 3 分
%% typeId: 1
%% type: 单选题
%% chapter: 原子和原子核
%% point: 半衰期
%% method: 无
%% score: 3
%% degree: 850
%% duplicateId: 0
%% body:
下列说法正确的是\xzanswer{D}
\fourchoices[answer=D]
{$\alpha$ 射线的穿透能力比 $\gamma$ 射线强}
{天然放射现象说明原子具有复杂的结构}
{核聚变中平均每个核子放出的能量比裂变中平均每个核子的小}
{半衰期跟放射性元素以单质或化合物形式存在无关}
%% answer: D
%% memo:
\memoanswer{
$\alpha$ 射线的穿透能力比 $\gamma$ 射线弱，故 A 错误；天然放射现象说明原子核具有复杂的结构，故 B 错误；核聚变中平均每个核子放出的能量比裂变中平均每个核子的大，故 C 错误；半衰期只与原子核的内部结构有关，与放射性元素以单质或化合物形式存在无关，故 D 正确。
}

\item
%% number: 5
%% paperName: 2020年1月浙江选考第 5 题 3 分
%% typeId: 1
%% type: 单选题
%% chapter: 曲线运动
%% point: 平抛运动
%% method: 无
%% score: 3
%% degree: 550
%% duplicateId: 0
%% body:
如图所示，钢球从斜槽轨道末端以 $v_{0}$ 的水平速度飞出，经过时间 $t$ 落在斜靠的挡板 AB 中点。若钢球以 $2v_{0}$ 的速度水平飞出，则\xzanswer{C}
\onepicture[w=5cm]{figs/05.png}
\fourchoices[answer=C]
{下落时间仍为 $t$}
{下落时间为 $2t$}
{下落时间为 $\sqrt 2 t$}
{落在挡板底端 B 点}
%% answer: C
%% memo:
\memoanswer{
设 A、B 两点的水平位移为 $2x$，则钢球以 $v_{0}$ 的速度水平抛出，下落时间为 $t$，水平位移为 $x=v_{0}t$；若钢球以 $2v_{0}$ 的速度水平抛出，下落时间 $t$ 时，其水平位移为 $x'=2v_{0}t=2x$，即 $t$ 时刻钢球在 B 点正上方，故钢球会落在挡板底端 B 点的右侧，故 D 错误。设 A、B 两点间竖直位移为 $H$，由 $h=\frac{1}{2}gt^{2}$ 得，钢球以 $v_{0}$ 的速度水平抛出的下落时间 $t=\sqrt{\frac{H}{g}}$，则钢球以 $2v_{0}$ 的速度水平抛出时下落时间 $t'=\sqrt{\frac{2H}{g}}=\sqrt{2}t$，故 C 正确，A、B 错误。
}

\item
%% number: 6
%% paperName: 2020年1月浙江选考第 6 题 3 分，即 2020年1月浙江学业水平考第 17 题 2 分
%% typeId: 1
%% type: 单选题
%% chapter: 电路
%% point: 电阻定律
%% method: 无
%% score: 3
%% degree: 750
%% duplicateId: 30655
%% body:
小明在一根细橡胶管中灌满食盐水，两端用粗铜丝塞住管口，形成一段封闭的盐水柱。他将此盐水柱接到电源两端，电源电动势和内阻恒定。握住盐水柱两端将它水平均匀拉伸到原长的 1.2 倍，若忽略温度对电阻率的影响，则此盐水柱\xzanswer{B}
\fourchoices[answer=B]
{通过的电流增大}
{两端的电压增大}
{阻值增大为原来的 1.2 倍}
{电功率增大为原来的 1.44 倍}
%% answer: B
%% memo:
\memoanswer{
将盐水柱拉长，盐水柱的电阻变大，由闭合电路欧姆定律知，通过盐水柱的电流减小，故 A 错误；由 $U=E-Ir$ 知，盐水柱两端的电压增大，故 B 正确；盐水柱的体积不变，将其水平均匀拉伸到原长的 1.2 倍，横截面积变为原来的 $\frac{1}{1.2}$，由 $R=\rho\frac{L}{S}$ 知，阻值增大为原来的 1.44 倍，故 C 错误；由 $P=I^{2}R$ 知，电流减小、电阻增大为原来的 1.44 倍，其电功率未增大为原来的 1.44 倍，故 D 错误。
}

\item
%% number: 7
%% paperName: 2020年1月浙江选考第 7 题 3 分
%% typeId: 1
%% type: 单选题
%% chapter: 电场
%% point: 带电粒子的偏转
%% method: 无
%% score: 3
%% degree: 550
%% duplicateId: 0
%% body:
如图所示，电子以某一初速度沿两块平行板的中线方向射入偏转电场中，已知极板长度 $l$，间距 $d$，电子质量 $m$，电荷量 $e$。若电子恰好从极板边缘射出电场，由以上条件可以求出的是\xzanswer{B}
\onepicture[w=5cm]{figs/07.png}
\fourchoices[answer=B]
{偏转电压}
{偏转的角度}
{射出电场速度}
{电场中运动的时间}
%% answer: B
%% memo:
\memoanswer{
设偏转电压为 $U$，初速度为 $v_{0}$，电子做类平抛运动，有 $l=v_{0}t$，$\frac{d}{2}=\frac{1}{2}at^{2}$，$e\frac{U}{d}=ma$，由以上各式解得 $U=\frac{md^{2}v_{0}^{2}}{el^{2}}$，由于初速度 $v_{0}$ 未知，故无法求出偏转电压 $U$，故 A 错误；$\tan\theta=\frac{v_{y}}{v_{0}}=\frac{at}{v_{0}}=\frac{d}{l}$，故可求出偏转角度，故 B 正确；由动能定理得 $e\frac{U}{2}=\frac{1}{2}mv^{2}-\frac{1}{2}mv_{0}^{2}$，解得 $v=\sqrt{1+\frac{d^{2}}{l^{2}}}\,v_{0}$，由于初速度 $v_{0}$ 未知，故无法求出电子射出电场的速度 $v$，故 C 错误；电子在电场中的运动时间 $t=\frac{l}{v_{0}}$，由于初速度 $v_{0}$ 未知，故无法求出 $t$，故 D 错误。
}

\item
%% number: 8
%% paperName: 2020年1月浙江选考第 8 题 3 分
%% typeId: 1
%% type: 单选题
%% chapter: 交流电
%% point: LC振荡回路
%% method: 无
%% score: 3
%% degree: 650
%% duplicateId: 0
%% body:
如图所示，单刀双掷开关 S 先打到 a 端让电容器充满电。$t=0$ 时开关 S 打到 b 端，$t=0.02\Us$ 时 LC 回路中电容器下极板带正电荷且电荷量第一次达到最大值。则\xzanswer{C}
\onepicture[w=5cm]{figs/08.png}
\fourchoices[answer=C]
{LC 回路的周期为 $0.02\Us$}
{LC 回路的电流最大时电容器中电场能最大}
{$t=1.01\Us$ 时线圈中磁场能最大}
{$t=1.01\Us$ 时回路中电流沿顺时针方向}
%% answer: C
%% memo:
\memoanswer{
开关 S 打到 a 端时电容器充电，电容器上极板带正电；开关打到 b 端后，$t=0.02\Us$ 时 LC 回路中电容器下极板带正电荷且电荷量第一次达到最大值，则该过程中电容器先放电后充电，为半个周期，故 LC 回路的周期 $T=2t=0.04\Us$，故 A 错误；LC 回路的电流最大时电容器中电场能为零，故 B 错误；$t=1.01\Us=25T+\frac{1}{4}T$，此时电容器放电完毕，线圈中磁场能最大，回路中电流方向沿逆时针方向，故 C 正确、D 错误。
}

\item
%% number: 9
%% paperName: 2020年1月浙江选考第 9 题 3 分
%% typeId: 1
%% type: 单选题
%% chapter: 曲线运动
%% point: 人造地球卫星
%% method: 无
%% score: 3
%% degree: 650
%% duplicateId: 0
%% body:
如图所示，卫星 a、b、c 沿圆形轨道绕地球运行。a 是极地轨道卫星，在地球两极上空约 1000 $\Ukm$ 处运行；b 是低轨道卫星，距地球表面高度与 a 相等；c 是地球同步卫星，则\xzanswer{C}
\onepicture[w=6cm]{\includesvg{figs/09}}
\fourchoices[answer=C]
{a、b 的周期比 c 大}
{a、b 的向心力一定相等}
{a、b 的速度大小相等}
{a、b 的向心加速度比 c 小}
%% answer: C
%% memo:
\memoanswer{
由万有引力定律得 $G\frac{Mm}{r^{2}}=m\frac{4\pi^{2}}{T^{2}}r$，解得卫星的运动周期 $T=2\pi\sqrt{\frac{r^{3}}{GM}}$，c 的轨道半径最大，故 c 的周期比 a、b 的周期大，故 A 错误；由万有引力定律知，向心力 $F=G\frac{Mm}{r^{2}}$，a、b 的质量大小关系未知，故向心力不一定相等，故 B 错误；由万有引力定律得 $G\frac{Mm}{r^{2}}=m\frac{v^{2}}{r}$，解得卫星的速度大小 $v=\sqrt{\frac{GM}{r}}$，a、b 的轨道半径相等，故速度大小相等，故 C 正确；由万有引力定律得 $G\frac{Mm}{r^{2}}=ma$，解得卫星的向心加速度大小 $a=\frac{GM}{r^{2}}$，c 的轨道半径最大，故 c 的向心加速度比 a、b 的小，故 D 错误。
}

\item
%% number: 10
%% paperName: 2020年1月浙江选考第 10 题 3 分
%% typeId: 1
%% type: 单选题
%% chapter: 交流电
%% point: 变压器
%% method: 无
%% score: 3
%% degree: 750
%% duplicateId: 0
%% body:
如图所示，甲乙两图中的理想变压器以不同的方式接在高压电路中。甲图中变压器原副线圈的匝数比为 $k_{1}$，电压表读数为 $U$，乙图中变压器原副线圈的匝数比为 $k_{2}$，电流表读数为 $I$。则甲图中高压线电压和乙图中高压线电流分别为\xzanswer{B}
\onepicture[w=7cm]{figs/10.png}
\fourchoices[answer=B]
{$k_{1}U$，$k_{2}I$}
{$k_{1}U$，$I/k_{2}$}
{$U/k_{1}$，$k_{2}I$}
{$U/k_{1}$，$I/k_{2}$}
%% answer: B
%% memo:
\memoanswer{
由 $\frac{U_{1}}{U_{2}}=\frac{n_{1}}{n_{2}}$ 得，题图甲中高压线电压为 $U'=k_{1}U$；由 $\frac{I_{1}}{I_{2}}=\frac{n_{2}}{n_{1}}$ 得，题图乙中高压线电流 $I'=\frac{I}{k_{2}}$，故 B 正确。
}

\item
%% number: 11
%% paperName: 2020年1月浙江选考第 11 题 3 分
%% typeId: 1
%% type: 单选题
%% chapter: 电磁感应
%% point: 楞次定律推广
%% method: 无
%% score: 3
%% degree: 550
%% duplicateId: 0
%% body:
如图所示，在光滑绝缘水平面上，两条固定的相互垂直彼此绝缘的导线通以大小相同的电流 $I$。在角平分线上，对称放置四个相同的正方形金属框。当电流在相同时间间隔内增加相同量，则\xzanswer{B}
\onepicture[w=5cm]{figs/11.png}
\fourchoices[answer=B]
{1、3 线圈静止不动，2、4 线圈沿着对角线向内运动}
{1、3 线圈静止不动，2、4 线圈沿着对角线向外运动}
{2、4 线圈静止不动，1、3 线圈沿着对角线向内运动}
{2、4 线圈静止不动，1、3 线圈沿着对角线向外运动}
%% answer: B
%% memo:
\memoanswer{
由安培定则及磁场的叠加知，穿过 1、3 线框的磁通量为零，穿过 2 线框的磁通量垂直纸面向外，穿过 4 线框的磁通量垂直纸面向里。当电流在相同时间内增加量相同时，由楞次定律知，1、3 线圈静止不动，2、4 线圈沿对角线向外运动，故 B 正确。
}

\item
%% number: 12
%% paperName: 2020年1月浙江选考第 12 题 3 分
%% typeId: 1
%% type: 单选题
%% chapter: 光学
%% point: 全反射
%% method: 无
%% score: 3
%% degree: 650
%% duplicateId: 0
%% body:
如图所示，一束光与某材料表面成 45${^\circ}$ 角入射，每次反射的光能量为入射光能量的 $k$ 倍（$0 < k < 1$）。若这束光最终进入材料的能量为入射光能量的（$1-k^{2}$）倍，则该材料折射率至少为\xzanswer{A}
\onepicture[w=5cm]{figs/12.png}
\fourchoices[answer=A]
{$\frac{\sqrt 6}{2}$}
{$\sqrt 2$}
{1.5}
{2}
%% answer: A
%% memo:
\memoanswer{
光最终进入材料的能量一定，故光最终发生全反射。光路图如图所示，根据折射定律得 $n=\frac{\sin45^{\circ}}{\sin\alpha}$，由 $\sin(90^{\circ}-\alpha)\geq\frac{1}{n}$，由以上两式解得 $n\geq\frac{\sqrt 6}{2}$，故 A 正确。

\onepicture[w=6cm, align=center]{figs/12_an.jpg}
}

\item
%% number: 13
%% paperName: 2020年1月浙江选考第 13 题 3 分
%% typeId: 1
%% type: 单选题
%% chapter: 电场
%% point: 带电粒子的平衡
%% method: 无
%% score: 3
%% degree: 450
%% duplicateId: 0
%% body:
如图所示，在倾角为 $\alpha$ 的光滑绝缘斜面上固定一个挡板，在挡板上连接一根劲度系数为 $k_{0}$ 的绝缘轻质弹簧，弹簧另一端与 A 球连接。A、B、C 三小球的质量均为 $M$，$q_{A}=q_{0}>0$，$q_{B}=-q_{0}$，当系统处于静止状态时，三小球等间距排列。已知静电力常量为 $k$，则\xzanswer{A}
\onepicture[w=6cm]{figs/13.png}
\fourchoices[answer=A]
{$q_{C}=\frac{4}{7}q_{0}$}
{弹簧伸长量为 $\frac{Mg\sin\alpha}{k_{0}}$}
{A 球受到的库仑力大小为 $2Mg$}
{相邻两小球间距为 $q_{0}\sqrt{\frac{3k}{7Mg}}$}
%% answer: A
%% memo:
\memoanswer{
设小球间的距离为 $r$，对 B 球，由平衡条件得 $k\frac{q_{0}^{2}}{r^{2}}=k\frac{q_{0}q_{C}}{r^{2}}+Mg\sin\alpha$；对 C 球，由平衡条件得 $k\frac{q_{0}q_{C}}{(2r)^{2}}+Mg\sin\alpha=k\frac{q_{0}q_{C}}{r^{2}}$，联立解得 $q_{C}=\frac{4}{7}q_{0}$，$r=q_{0}\sqrt{\frac{3k}{7Mg\sin\alpha}}$，故 A 正确、D 错误；A、B、C 球作为整体，由平衡条件得 $k_{0}x=3Mg\sin\alpha$，解得弹簧的伸长量 $x=\frac{3Mg\sin\alpha}{k_{0}}$，故 B 错误；对 A 球，由平衡条件得 $k_{0}x=Mg\sin\alpha+F$，解得 A 球受到的库仑力大小 $F=2Mg\sin\alpha$，故 C 错误。
}

\item
%% number: 14
%% paperName: 2020年1月浙江选考第 14 题 2 分
%% typeId: 2
%% type: 多选题
%% chapter: 原子和原子核
%% point: 玻尔模型
%% method: 无
%% score: 2
%% degree: 750
%% duplicateId: 0
%% body:
由玻尔原子模型求得氢原子能级如图所示，已知可见光的光子能量在 1.62 $\UeV$ 到 3.11 $\UeV$ 之间，则\xzanswer{CD}
\onepicture[w=4.5cm]{\includesvg{figs/14}}
\fourchoices[answer=CD]
{氢原子从高能级向低能级跃迁时可能辐射出 $\gamma$ 射线}
{氢原子从 $n=3$ 的能级向 $n=2$ 的能级跃迁时会辐射出红外线}
{处于 $n=3$ 能级的氢原子可以吸收任意频率的紫外线并发生电离}
{大量氢原子从 $n=4$ 能级向低能级跃迁时可辐射出 2 种频率的可见光}
%% answer: CD
%% memo:
\memoanswer{
$\gamma$ 射线是原子核受激发产生的，氢原子跃迁不能辐射出 $\gamma$ 射线，故 A 错误；氢原子从 $n=3$ 能级向 $n=2$ 能级跃迁时辐射出光子的能量为 $\Delta E=-1.51\UeV-(-3.4\UeV)=1.89\UeV$，而可见光的光子能量在 1.62 $\UeV$ 到 3.11 $\UeV$ 之间，故该跃迁不会辐射出红外线，故 B 错误；处于 $n=3$ 能级的氢原子的电离能为 1.51 $\UeV$，紫外线的能量大于 3.11 $\UeV$，故该氢原子可吸收任意频率的紫外线并发生电离，故 C 正确；从 $n=4$ 能级到 $n=3$ 能级辐射光子的能量 $\Delta E_{1}=-0.85\UeV-(-1.51\UeV)=0.66\UeV$，从 $n=4$ 能级到 $n=2$ 能级辐射光子的能量 $\Delta E_{2}=-0.85\UeV-(-3.4\UeV)=2.55\UeV$，从 $n=4$ 能级到 $n=1$ 能级辐射光子的能量 $\Delta E_{3}=-0.85\UeV-(-13.6\UeV)=12.75\UeV$，从 $n=3$ 能级到 $n=2$ 能级辐射光子的能量 $\Delta E_{4}=-1.51\UeV-(-3.4\UeV)=1.89\UeV$，从 $n=3$ 能级到 $n=1$ 能级辐射光子的能量 $\Delta E_{5}=-1.51\UeV-(-13.6\UeV)=12.09\UeV$，从 $n=2$ 能级到 $n=1$ 能级辐射光子的能量 $\Delta E_{6}=-3.4\UeV-(-13.6\UeV)=10.2\UeV$，故可辐射出 2 种频率的可见光，故 D 正确。
}

\item
%% number: 15
%% paperName: 2020年1月浙江选考第 15 题 2 分
%% typeId: 2
%% type: 多选题
%% chapter: 光学
%% point: 光的折射
%% method: 无
%% score: 2
%% degree: 550
%% duplicateId: 0
%% body:
如图所示，波长为 $\lambda_{a}$ 和 $\lambda_{b}$ 的两种单色光射入三棱镜，经折射后射出两束单色光 a 和 b，则这两束光\xzanswer{BD}
\onepicture[w=5cm]{\includesvg{figs/15}}
\fourchoices[answer=BD]
{照射同一种金属均有光电子逸出，光电子最大初动能 $E_{ka} > E_{kb}$}
{射向同一双缝干涉装置，其干涉条纹间距 $\Delta x_{a} > \Delta x_{b}$}
{在水中的传播速度 $v_{a} < v_{b}$}
{光子动量 $p_{a} < p_{b}$}
%% answer: BD
%% memo:
\memoanswer{
由题图知，三棱镜对 a 光的偏折程度小，故三棱镜对 a 光的折射率小，a 光的频率小，由光电效应方程 $E_{\text{km}}=h\nu-W_{\text{逸}}$ 知，照射同一种金属时逸出的光电子最大初动能 $E_{ka}<E_{kb}$，故 A 错误；a 光的频率小、波长长，由 $\Delta x=\frac{L}{d}\lambda$ 知，射向同一双缝干涉装置时其干涉条纹间距 $\Delta x_{a}>\Delta x_{b}$，故 B 正确；由 $n=\frac{c}{v}$ 知，在水中的传播速度 $v_{a}>v_{b}$，故 C 错误；由 $p=\frac{h}{\lambda}$ 知，光子动量 $p_{a}<p_{b}$，故 D 正确。
}

\item
%% number: 16
%% paperName: 2020年1月浙江选考第 16 题 2 分
%% typeId: 2
%% type: 多选题
%% chapter: 机械波
%% point: 波长频率波速
%% method: 无
%% score: 2
%% degree: 550
%% duplicateId: 0
%% body:
如图所示，波源 O 垂直于纸面做简谐运动，所激发的横波在均匀介质中向四周传播，图中虚线表示两个波面。$t=0$ 时，离 O 点 5 $\Um$ 的 A 点开始振动；$t=1\Us$ 时，离 O 点 10 $\Um$ 的 B 点也开始振动，此时 A 点第五次回到平衡位置，则\xzanswer{AB}
\onepicture[w=5cm]{\includesvg{figs/16}}
\fourchoices[answer=AB]
{波的周期为 $0.4\Us$}
{波的波长为 $2\Um$}
{波速为 $5\sqrt 3\Ums$}
{$t=1\Us$ 时 AB 连线上有 4 个点处于最大位移}
%% answer: AB
%% memo:
\memoanswer{
$t=1\Us$ 时 A 点第五次回到平衡位置，则 $\frac{5}{2}T=1\Us$，解得 $T=0.4\Us$，故 A 正确；波的传播速度 $v=\frac{\overline{OB}-\overline{OA}}{t}=5\Ums$，故 C 错误；波长 $\lambda=vT=2\Um$，故 B 正确；因为题中未明确给出 A、B 两点具体位置关系，故 AB 间处于最大位移处的质点个数不确定，但一定不少于 5 个，故 D 错误。
}

\item
%% number: 17
%% paperName: 2020年1月浙江选考第 17 题 7 分
%% typeId: 4
%% type: 实验题
%% chapter: 运动和力
%% point: 验证牛顿第二定律
%% method: 无
%% score: 7
%% degree: 650
%% duplicateId: 0
%% body:
在“探究加速度与力、质量的关系”和用橡皮筋“探究做功与物体速度变化的关系”实验中
\begin{enumerate}
	\item
	都是通过分析纸带上的点来测量物理量，下列说法正确的是\_\_\_\_\_\_\_\_（多选）

	A．都需要分析打点计时器打下的第一个点

	B．都不需要分析打点计时器打下的第一个点

	C．一条纸带都只能获得一组数据

	D．一条纸带都能获得多组数据

	\item
	如图是两条纸带的一部分，A、B、C、$\ldots$、G 是纸带上标出的计数点，每两个相邻的计数点之间还有 4 个打出的点未画出。其中图\_\_\_\_（填“甲”或“乙”）所示的是用橡皮筋“探究做功与物体速度变化的关系”的实验纸带。“探究加速度与力、质量的关系”实验中，小车的加速度大小 $a=$\_\_\_\_$\Umsq$（保留 2 位有效数字）。
	\item
	在用橡皮筋“探究做功与物体速度变化的关系”实验中，平衡阻力后，小车与橡皮筋组成的系统在橡皮筋恢复形变前机械能\_\_\_\_（填“守恒”或“不守恒”）。
\end{enumerate}
\twopicture[numstyle=chinese, wA=11cm, wB=11cm, align=center]{figs/17a.png}{figs/17b.png}
%% answer: 
\jdanswer{
\begin{enumerate}
	\item
	BC

	\item
	甲，0.40

	\item
	不守恒

\end{enumerate}
}
%% memo:
\memoanswer{
（1）“探究加速度与力、质量的关系”实验中利用纸带求小车的加速度，“探究做功与物体速度变化的关系”实验中利用纸带求小车的最大速度，故都不需要分析打点计时器打下的第一个点，故 B 正确、A 错误；一条纸带都只能获得一组数据，故 C 正确、D 错误。

（2）“探究做功与物体速度变化的关系”实验中，橡皮筋恢复原长后，小车做匀速直线运动，故题图甲是用橡皮筋“探究做功与物体速度变化的关系”的实验纸带；相邻计数点间还有 4 个点未画出，则相邻计数点间的时间间隔 $t=0.1\Us$，由题图乙知，小车的加速度 $a=\frac{DE+EF+FG-AB-BC-CD}{(3T)^{2}}=\frac{0.0220+0.0260+0.0300-0.0100-0.0140-0.0179}{(3\times0.1)^{2}}\Umsq\approx0.40\Umsq$。

（3）平衡阻力后，小车和橡皮筋组成的系统在橡皮筋恢复形变前受阻力作用，且阻力做功，系统机械能不守恒。
}

\item
%% number: 18
%% paperName: 2020年1月浙江选考第 18 题 7 分
%% typeId: 4
%% type: 实验题
%% chapter: 电路
%% point: 伏安法测电阻
%% method: 无
%% score: 7
%% degree: 550
%% duplicateId: 0
%% body:
（1）小明同学用多用电表测量一未知电阻器的阻值。经过规范操作后，所选欧姆挡倍率及指针位置分别如图甲、乙所示，则此电阻器的阻值为\_\_\_\_$\UO$。

在“测绘小灯泡的伏安特性曲线”实验中：

①如图丙所示，已经连接了一部分电路，请在答题纸上对应位置将电路连接完整。

②合上开关后，测出 9 组 $I$、$U$ 值，在 $I-U$ 坐标系中描出各对应点，如图丁所示。请在答题纸对应位置的图中画出此小灯泡的伏安特性曲线。

③与图丁中 P 点对应的状态，小灯泡灯丝阻值最接近\_\_\_\_

（A）$16.7\UO$　　（B）$12.4\UO$　　（C）$6.2\UO$

\fourpicture[numstyle=chinese, wA=4.6cm, wB=4.6cm, wC=4.6cm, wD=4.6cm, align=center]{figs/18a.png}{figs/18b.png}{figs/18c.png}{figs/18d.png}
%% answer: 
\jdanswer{
\begin{enumerate}
	\item
	1700$\sim$1800

	\item
	①连线正确；②作图正确；③C
\end{enumerate}
}
%% memo:
\memoanswer{
（1）由题图甲、乙知，电阻器的阻值约为 $R=17.5\times100\UO=1750\UO$。

（2）①连线如答图甲所示；②作图如答图乙所示；③由题图丁知，P 点对应电压约为 $U=1.40\UV$，电流约为 $I=0.23\UA$，由欧姆定律得 $R=\frac{U}{I}\approx6.1\UO$，故 C 正确。

\onepicture[w=10cm, align=center]{figs/18_an.jpg}
}

\item
%% number: 20
%% paperName: 2020年1月浙江选考第 20 题 9 分
%% typeId: 6
%% type: 计算题
%% chapter: 运动和力
%% point: 加速减速模型
%% method: 无
%% score: 9
%% degree: 550
%% duplicateId: 0
%% body:
一个无风晴朗的冬日，小明乘坐游戏滑雪车从静止开始沿斜直雪道下滑，滑行 54 $\Um$ 后进入水平雪道，继续滑行 40.5 $\Um$ 后减速到零。已知小明和滑雪车的总质量为 60 $\Ukg$，整个滑行过程用时 10.5 $\Us$，斜直雪道倾角为 37${^\circ}$（$\sin37{^\circ}=0.6$）。求小明和滑雪车
\begin{enumerate}
	\item
	滑行过程中的最大速度 $v_{m}$ 的大小；
	\item
	在斜直雪道上滑行的时间 $t_{1}$；
	\item
	在斜直雪道上受到的平均阻力 $F_{f}$ 的大小。
\end{enumerate}
\onepicture[w=7cm, align=right]{figs/20a.jpg}
%% answer: 
\jdanswer{
\begin{enumerate}
	\item
	$v_{m}=18\Ums$

	\item
	$t_{1}=6\Us$

	\item
	$F_{f}=180\UN$

\end{enumerate}
}
%% memo:
\memoanswer{
（1）由题意知，滑雪车先做匀加速运动后做匀减速运动，加速阶段与减速阶段的平均速度相同。设滑行过程中的最大速度为 $v_{m}$，则平均速度为 $\frac{v_{m}}{2}$，由运动学公式得 $\frac{v_{m}}{2}t=x_{1}+x_{2}$，解得 $v_{m}=18\Ums$。

（2）同（1），由运动学公式得 $x_{1}=\frac{v_{m}}{2}t_{1}$，解得 $t_{1}=6\Us$。

（3）设滑雪车在斜直雪道上运动时的加速度为 $a$，则 $a=\frac{v_{m}}{t_{1}}=3\Umsq$，由牛顿第二定律得 $mg\sin37^{\circ}-F_{f}=ma$，解得 $F_{f}=180\UN$。
}

\item
%% number: 20
%% paperName: 2020年1月浙江选考第 20 题 12 分，即 2020年1月浙江学业水平考第 23 题 9 分
%% typeId: 6
%% type: 计算题
%% chapter: 功和能
%% point: 机械能守恒定律
%% method: 无
%% score: 12
%% degree: 550
%% duplicateId: 30661
%% body:
如图所示，一弹射游戏装置由安装在水平台面上的固定弹射器、竖直圆轨道（在最低点 E 分别与水平轨道 EO 和 EA 相连）、高度 $h$ 可调的斜轨道 AB 组成。游戏时滑块从 O 点弹出，经过圆轨道并滑上斜轨道。全程不脱离轨道且恰好停在 B 端则视为游戏成功。已知圆轨道半径 $r=0.1\Um$，OE 长 $L_{1}=0.2\Um$，AC 长 $L_{2}=0.4\Um$，圆轨道和 AE 光滑，滑块与 AB、OE 之间的动摩擦因数 $\mu=0.5$。滑块质量 $m=2\Ug$ 且可视为质点，弹射时从静止释放且弹簧的弹性势能完全转化为滑块动能。忽略空气阻力，各部分平滑连接。求：
\begin{enumerate}
	\item
	滑块恰好能过圆轨道最高点 F 时的速度 $v_{F}$ 大小；
	\item
	当 $h=0.1\Um$ 且游戏成功时，滑块经过 E 点对圆轨道的压力 $F_{N}$ 大小及弹簧弹性势能 $E_{p0}$；
	\item
	要使游戏成功，弹簧的弹性势能 $E_{p}$ 与高度 $h$ 之间满足的关系。
\end{enumerate}
\onepicture[w=8cm, align=right]{figs/20b.png}
%% answer: 
\jdanswer{
\begin{enumerate}
	\item
	$v_{F}=1\Ums$

	\item
	$F_{N}=0.14\UN$，$E_{p0}=8.0\times10^{-3}\UJ$

	\item
	$E_{p}=2\times10^{-3}(10h+3)\UJ$，其中 $0.05\Um\leq h\leq0.2\Um$

\end{enumerate}
}
%% memo:
\memoanswer{
（1）滑块恰好经过最高点 F 的条件为重力提供向心力，由牛顿第二定律得 $mg=m\frac{v_{F}^{2}}{r}$，解得 $v_{F}=1\Ums$。

（2）滑块从 E 到 B，由动能定理得 $-mgh-\mu mgL_{2}=0-\frac{1}{2}mv_{E}^{2}$；设滑块在 E 点受到的支持力为 $F_{N}$，由牛顿第二定律得 $F_{N}-mg=m\frac{v_{E}^{2}}{r}$，联立解得 $F_{N}=0.14\UN$。滑块从 O 到 B 点，由动能定理得 $E_{p0}-mgh-\mu mg(L_{1}+L_{2})=0$，解得 $E_{p0}=8.0\times10^{-3}\UJ$。

（3）由功能关系得，滑块恰能过 F 点的弹性势能为 $E_{p1}=2mgr+\mu mgL_{1}+\frac{1}{2}mv_{F}^{2}=7.0\times10^{-3}\UJ$；滑块运动到 B 点减速到 0，由功能关系得 $E_{p1}-mgh_{1}-\mu mg(L_{1}+L_{2})=0$，解得 $h_{1}=0.05\Um$。滑块恰好能停在 B 点，则 $\mu mg\cos\theta=mg\sin\theta$，解得 $\tan\theta=0.5$，此时 $h_{2}=0.2\Um$。滑块从 O 到 B 点，由功能关系得 $E_{p}=mgh+\mu mg(L_{1}+L_{2})=2\times10^{-3}(10h+3)\UJ$，其中 $0.05\Um\leq h\leq0.2\Um$。
}

\item
%% number: 21
%% paperName: 2020年1月浙江选考第 21 题 10 分
%% typeId: 6
%% type: 计算题
%% chapter: 电磁感应
%% point: 电磁感应中的变加速
%% method: 无
%% score: 10
%% degree: 450
%% duplicateId: 0
%% body:
如图甲所示，在 $xOy$ 水平面内，固定放置着间距为 $l$ 的两平行金属直导轨，其间连接有阻值为 $R$ 的电阻，电阻两端连接示波器（内阻可视为无穷大），可动态显示电阻 $R$ 两端的电压。两导轨间存在大小为 $B$、方向垂直导轨平面的匀强磁场。$t=0$ 时一质量为 $m$、长为 $l$ 的导体棒在外力 $F$ 作用下从 $x=x_{0}$ 位置开始做简谐运动，观察到示波器显示的电压随时间变化的波形是如图乙所示的正弦曲线。取 $x_{0}=-\frac{U_{m}T}{2\pi Bl}$，则简谐运动的平衡位置在坐标原点 O。不计摩擦阻力和其它电阻，导体棒始终垂直导轨运动。（提示：可以用 $F-x$ 图象下的“面积”代表力 $F$ 所做的功）
\begin{enumerate}
	\item
	求导体棒所受到的安培力 $F_{A}$ 随时间 $t$ 的变化规律；
	\item
	求在 0 至 0.25$T$ 时间内外力 $F$ 的冲量；
	\item
	若 $t=0$ 时外力 $F_{0}=1\UN$，$l=1\Um$，$T=2\pi\Us$，$m=1\Ukg$，$R=1\UO$，$U_{m}=0.5\UV$，$B=0.5\UT$，求外力与安培力大小相等时棒的位置坐标和速度。
\end{enumerate}
\twopicture[numstyle=chinese, wA=6.5cm, wB=6.5cm, align=center]{figs/21a.png}{figs/21b.png}
%% answer: 
\jdanswer{
\begin{enumerate}
	\item
	$F_{A}=-\frac{BlU_{m}}{R}\sin\frac{2\pi}{T}t$

	\item
	$I_{F}=\frac{BlU_{m}T}{2\pi R}+\frac{mU_{m}}{Bl}$

	\item
	$x_{1}'=\frac{1}{\sqrt 5}\Um$，$v_{1}'=\frac{2}{\sqrt 5}\Ums$ 和 $x_{2}'=-\frac{1}{\sqrt 5}\Um$，$v_{2}'=-\frac{2}{\sqrt 5}\Ums$

\end{enumerate}
}
%% memo:
\memoanswer{
（1）由显示的波形可得 $U=U_{m}\sin\frac{2\pi}{T}t$，回路中的电流大小为 $I=\frac{U_{m}}{R}\sin\frac{2\pi}{T}t$，导体棒受到的安培力大小为 $F_{A}=-BIl=-\frac{BlU_{m}}{R}\sin\frac{2\pi}{T}t$。

（2）由安培力产生过程可得安培力的冲量为 $I_{A}=Bl\Delta q=-\frac{B^{2}|x_{0}|l^{2}}{R}$；由（1）知 $U=U_{m}\sin\frac{2\pi}{T}t$，由法拉第电磁感应定律得 $U=Blv$，解得 $v=\frac{U}{Bl}$，即 $v=\frac{U_{m}}{Bl}\sin\frac{2\pi}{T}t$；从 0 至 0.25$T$，由动量定理得 $I_{F}+I_{A}=mv_{m}$，代入解得 $I_{F}=\frac{BlU_{m}T}{2\pi R}+\frac{mU_{m}}{Bl}$。

（3）由题意知，导体棒做简谐运动，则有 $F_{\text{回复}}=F_{A}+F=-kx$①。由题给条件知 $F_{0}=-kx_{0}$，解得 $k=1$。当 $F_{A}=-F$ 时，代入①可知 $x=0$，此时导体棒处于平衡位置，导体棒的速度取最大值，即 $v=\pm v_{m}=\pm1\Ums$；当 $F_{A}=F$ 时，设此时导体棒的位移为 $x'$、速度为 $v'$，代入①有 $F_{A}=-\frac{1}{2}kx'=-\frac{1}{2}x'$，又 $F_{A}=-BIl=-\frac{B^{2}l^{2}v}{R}=-\frac{1}{4}v$，联立化简得 $2x'=v'$②。设导体棒从 $x_{0}$ 运动至 $x'$ 时，回复力对导体棒做的功为 $W$，$F_{\text{回复}}$ 与 $x$ 成线性关系，由提示知 $F-x$ 图象下的“面积”代表 $F$ 所做的功（也可理解成平均力做功），则 $W=\frac{kx_{0}+kx'}{2}(x_{0}-x')=\frac{1}{2}k(x_{0}^{2}-x'^{2})$；导体棒从 $x_{0}$ 运动至 $x'$ 时，由动能定理得 $\frac{1}{2}mv'^{2}=\frac{1}{2}k(x_{0}^{2}-x'^{2})$③，联立②③解得 $x_{1}'=\frac{1}{\sqrt 5}\Um$ 或 $x_{2}'=-\frac{1}{\sqrt 5}\Um$，当 $x_{1}'=\frac{1}{\sqrt 5}\Um$ 时，$v_{1}'=\frac{2}{\sqrt 5}\Ums$；当 $x_{2}'=-\frac{1}{\sqrt 5}\Um$ 时，$v_{2}'=-\frac{2}{\sqrt 5}\Ums$。
}

\item
%% number: 22
%% paperName: 2020年1月浙江选考第 22 题 10 分
%% typeId: 6
%% type: 计算题
%% chapter: 磁场
%% point: 带电粒子在磁场中的运动
%% method: 无
%% score: 10
%% degree: 450
%% duplicateId: 0
%% body:
通过测量质子在磁场中的运动轨迹和打到探测板上的计数率（即打到探测板上质子数与衰变产生总质子数 $N$ 的比值），可研究中子（$\ce{^{1}_{0}n}$）的 $\beta$ 衰变。中子衰变后转化成质子和电子，同时放出质量可视为零的反中微子 $\overline{\nu}$。如图所示，位于 P 点的静止中子经衰变可形成一个质子源，该质子源在纸面内各向均匀地发射 $N$ 个质子。在 P 点下方放置有长度 $L=1.2\Um$ 以 O 为中点的探测板，P 点离探测板的垂直距离 OP 为 $a$。在探测板的上方存在方向垂直纸面向里、磁感应强度大小为 $B$ 的匀强磁场。

已知电子质量 $m_{e}=9.1\times10^{-31}\Ukg=0.51\UMeVcSq$，中子质量 $m_{n}=939.57\UMeVcSq$，质子质量 $m_{p}=938.27\UMeVcSq$（$c$ 为光速，不考虑粒子之间的相互作用）。

若质子的动量 $p=4.8\times10^{-21}\Ukg\cdot\Um\cdot\Us^{-1}=3\times10^{-8}\UMeV\cdot\Us\cdot\Um^{-1}$。
\begin{enumerate}
	\item
	写出中子衰变的核反应式，求电子和反中微子的总动能（以 $\UMeV$ 为能量单位）；
	\item
	当 $a=0.15\Um$，$B=0.1\UT$ 时，求计数率；
	\item
	若 $a$ 取不同的值，可通过调节 $B$ 的大小获得与（2）问中同样的计数率，求 $B$ 与 $a$ 的关系并给出 $B$ 的取值范围。
\end{enumerate}
\onepicture[w=7cm, align=right]{figs/22.png}
%% answer: 
\jdanswer{
\begin{enumerate}
	\item
	$\ce{^{1}_{0}n}\rightarrow\ce{^{1}_{1}p}+\ce{^{0}_{-1}e}+\overline{\nu}_{e}$，$E_{e}+E_{\nu}=0.7468\UMeV$

	\item
	$\eta=\frac{2}{3}$

	\item
	$B=\frac{3}{200a}\UT$，$B\geq\frac{\sqrt{15}}{40}\UT$
\end{enumerate}
}
%% memo:
\memoanswer{
（1）由题意知，中子衰变的核反应方程为 $\ce{^{1}_{0}n}\rightarrow\ce{^{1}_{1}p}+\ce{^{0}_{-1}e}+\overline{\nu}_{e}$。中子衰变时的质量亏损为 $\Delta m=m_{n}-(m_{p}+m_{e})$，由质能方程得 $\Delta E_{d}=[m_{n}-(m_{p}+m_{e})]c^{2}=0.79\UMeV$。由动量与动能的关系知，质子的动能为 $E_{kp}=\frac{p^{2}}{2m_{p}}=0.0432\UMeV$。由能量守恒得 $E_{e}+E_{\overline{\nu}_{e}}=\Delta E_{d}-E_{kp}=0.7468\UMeV$。

（2）质子运动时，由洛伦兹力提供向心力得质子运动的轨迹半径为 $qvB=m\frac{v^{2}}{R}$，解得 $R=\frac{mv}{qB}=\frac{p}{eB}=0.3\Um$，由题意知此时满足 $R=2a$；如图甲所示，质子打到探测板的两种临界情况对应发射角度分别为 $\alpha=\beta=\frac{\pi}{6}$，利用打到探测板与未打到探测板上的粒子的发射角度，可得质子计数率为 $\eta=\frac{\frac{4\pi}{3}}{2\pi}=\frac{2}{3}$。

（3）为确保计数率仍为 $\eta=\frac{2}{3}$，则仍需满足 $R'=2a$，即 $R'=\frac{p}{eB'}=2a$，解得 $B=\frac{3}{200a}\UT$。将图甲中的大圆轨迹绕 P 点逆时针旋转，若探测板足够长，轨迹均可以与探测板相交，即质子可以打到探测板上，当交点与 P 的连线长度等于直径时为临界情况，如图乙所示，恰能打到探测板左端的条件由几何关系得 $4R_{\max}^{2}-a^{2}=\frac{L^{2}}{4}$，又 $R_{\max}=2a$，联立解得 $B\geq\frac{\sqrt{15}}{40}\UT$。

\twopicture[numstyle=chinese, wA=6cm, wB=6cm, align=center]{figs/22a_an.png}{figs/22b_an.png}
}

\end{enumerate}

\end{document}
